% Pursuing Stacks (cotes 134-2 to 134-8): archive cote 134-7, site page 105 (PDF page 106), typescript page 547.
% typescript. Draft compared: ps6.tex:5272-5351.
% AI-transcribed by Claude Opus 5.5 (claude-opus-5-5, adaptive thinking, effort high), Message Batches API, batch ps4 (msgbatch_012Aoc1NuskATPsyBeKSYK7J), 2026-11-12; not yet checked by eye.
\scanpage{134-7/105}
\mspage{547}

I would like however to suggest still another extension of the noti-%
on of a "bundle", which maybe will prove something better than just
a random way out of embarassment ! It has to do with an attempt to come
to a description of one among the four "basic functors" of section 124,
namely the "spherical" functor
\[ (1) \qquad \overset{\hand{\sim}}{S} : \ \mathrm{Hotab}_{o}(k) \hand{\longrightarrow} \mathrm{Hot}_{o}(k) \quad , \]
which for the time being is remaining hypothetical, due notably to my
unability so far to carry out the suspension operation for schematic
homotopy types. There are two basic formal properties giving us some
clues about this functor, namely it should be left adjoint to the
\hand{``}total homotopy\hand{''} or \hand{``}Lie\hand{''} functor $L\hand{\Pi}.$ , and it should be right inverse to
the \inserted{(reduced)} total homology functor $L\overset{\hand{\sim}}{H}.$ (cf. \over{\unsure{b}}{p}ages 523-525). \unsure{Let's|Lat's} work for
the sake of definiteness, \xout{with} for describing $\mathrm{Hot}_{o}(k)$, with semisimpli-%
cial V-bundles $X_{*}$ satisfying the extra assumption $X_{o}$=e (and presumably
a little more), corresponding therefore to co-semisimplicial algebras
$A^{*}$ satisfying $A^{o} = k$ , such that the components $A^{n}$ be isomorphic to
symmetric algebras
\[ (2) \qquad A^{n} \hand{\simeq} \mathrm{Sym}\xout{*}_{k}(M^{n}) \ , \]
where we'll assume the modules $M^{n}$ (or equivalently, the algebras $A^{n}$)
\uline{flat} over k. The condition $A^{o}$=k implies that the algebras $A^{n}$ are k-aug-%
mented, and i\over{\unsure{s}}{f} $\underline{m}^{\hand{(n)}}_{\deleted{*}}$ is the augmentation ideal, we get
\[ (3) \qquad \underline{m}^{\hand{(n)}}/(\underline{m}^{\hand{(n)}})^{\hand{2}} \hand{\simeq} M^{n} \quad , \]
thus the modules $M^{n}$ may be viewed as the components of a co-semisimplicial
k-module $M^{*}$ , giving rise to a semisimplicial vector bundle $V(M^{*})$ over
k , and we may write
\[ (4) \qquad L\hand{\Pi}.(X_{*}) \hand{\simeq} V(M^{*}) \quad , \]
where \inserted{\deleted{\illegible{1}}} the second member is viewed as a chain complex (rather than as a
semisimplicial Module) in the usual way. \xout{Let now} Consider now an object
\xout{\unsure{an}} $L_{*}$ in Hotab$_{o}$(k), it is described in terms of a \inserted{co-}semisimplicial
k-module $L^{*}$ by the formula
\[ (5) \qquad L_{*} = V(L^{*}) \quad , \]
and we get now
\[ (\hand{*}) \qquad \mathrm{Hom}_{\underline{O}_{k}}(L_{*},L\hand{\Pi}.(X_{*})) \hand{\simeq} \mathrm{Hom}_{\underline{O}_{k}}(V(L^{*}),V(M^{*})) \hand{\simeq} \mathrm{Hom}_{\underline{O}_{k}}(M^{*},L^{*}) \quad . \]
Now if L is any k-module, let
\[ (6) \qquad D(M) = \mathrm{Sym}_{k}(M)(1) \hand{\simeq} k \oplus M \]

\margin{105}
